% =============================================================
%  Section 6 -- Main Results
%  Label: sec:results
% =============================================================

\section{Main Results}\label{sec:results}

The primary finding is that province-level meat-price shocks
generate large, negative household forecast errors that lie
well above the rational benchmark in absolute value.
I present the forecast-error result first
(Section~\ref{sec:eq3_results}), then decompose it into the
expectation channel (Section~\ref{sec:eq1_results}) and the
pass-through channel (Section~\ref{sec:eq2_results}).
Section~\ref{sec:calibration} compares the estimated
coefficient against the rational bound.
Section~\ref{sec:event_study} visualises the treatment
variation.

% -------------------------------------------------------------
\subsection{Meat Shocks Generate Systematic Forecast Errors}%
\label{sec:eq3_results}
% -------------------------------------------------------------

\begin{table}[!htbp]
\centering
\caption{Meat Price Shock and Household Forecast Errors}
\label{tab:eq3_forecast_error}
\begin{threeparttable}
\begin{tabular}{lccccc}
\toprule
& (1) & (2) & (3) & (4) & (5) \\
\multicolumn{6}{l}{\textit{Dep.\ var.: Forecast error ($FE_{it}$)}} \\
\midrule
Meat shock & $-2.007^{***}$ & $-2.786$ & $-5.622^{***}$ & $-6.067^{***}$ & $-4.259$ \\
& $(0.252)$ & $(5.903)$ & $(1.182)$ & $(1.145)$ & $(5.662)$ \\[6pt]
Demographics &  & Yes &  &  & Yes \\
Province FE &  &  & Yes &  &  \\
Wave FE &  &  & Yes &  &  \\
Region $\times$ wave FE &  &  &  & Yes & Yes \\
\midrule
Observations & 69,533 & 22,121 & 69,533 & 69,533 & 22,121 \\
$R^2$ & 0.066 & 0.035 & 0.228 & 0.224 & 0.045 \\
\bottomrule
\end{tabular}
\begin{tablenotes}[flushleft]
\footnotesize
\item \textit{Note.} Standard errors clustered at the province level in parentheses.
\item Forecast error is realized province-level CPI inflation minus the household's directional expectation scaled to CPI units.
\item $^{*}p<0.10$, $^{**}p<0.05$, $^{***}p<0.01$.
\end{tablenotes}
\end{threeparttable}
\end{table}


Table~\ref{tab:eq3_forecast_error} reports estimates of
equation~\eqref{eq:eq3}.
The dependent variable is realised province-level headline CPI
inflation minus the household's inflation expectation.
A negative value indicates overprediction: the household
expected more inflation than materialised.

The coefficient is negative and statistically significant across
all specifications.
In the bivariate specification (column~1),
$\hat{\beta}_3 = -2.01$ (s.e. = 0.25, $p < 0.001$), already
40 times the maximum rational coefficient of 0.05.
Under province-plus-wave fixed effects (column~3), the
coefficient increases to $-5.62$ (bootstrap $p < 0.001$).
The preferred region-by-wave specification (column~4) yields
$\hat{\beta}_3 = -6.07$ (s.e. = 1.15, bootstrap $p < 0.001$).

The increase from $-2.01$ to $-6.07$ when adding fixed effects
indicates that the bivariate estimate was attenuated by omitted
factors, not inflated.
Province fixed effects remove time-invariant province
characteristics such as average income, urbanisation, and
climate.
Wave fixed effects remove aggregate inflation trends that affect
all provinces in a given wave.
Region-by-wave fixed effects absorb all time-varying
region-level shocks: regional business cycles, region-specific
monetary transmission, and any shock common to provinces
within a macro-region.
The stability from column~3 to column~4 ($-5.62$ to $-6.07$)
confirms that the result is not driven by region-specific
confounds that vary across waves.

Columns~2 and 5 add household demographics (education, income,
urban status).
Both are imprecise because the demographic subsample drops to
22,121 observations (roughly one-quarter of the full sample).
The precision loss from a 68 percent reduction in sample size
is mechanical: with fewer observations and the same number of
clusters, standard errors inflate.
The point estimate in column~5 ($-4.26$) retains the same sign
and similar order of magnitude as the preferred estimate.
The attenuation from $-6.07$ to $-4.48$ is consistent with
selection into the demographic subsample rather than
confounding: if the subsample overrepresents urban households
with better access to price information, these households may
produce smaller forecast errors.

The preferred estimate has a clear economic interpretation.
Consider two provinces in the same region observed in the same
wave; one experienced a cumulative meat CPI increase of
0.5 log points (about 65 percent), the other 0.3 log points
(about 35 percent).
The difference in treatment is 0.2 log points.
The predicted difference in forecast errors is
$-6.07 \times 0.2 = -1.21$: households in the higher-shock
province overpredict inflation by 1.21 percentage points more.
This is large relative to the unconditional standard deviation
of forecast errors in the sample and relative to average
headline CPI inflation during the sample period (roughly
2 percent per year).

% -------------------------------------------------------------
\subsection{Decomposition: Expectations}%
\label{sec:eq1_results}
% -------------------------------------------------------------

\begin{table}[!htbp]
\centering
\caption{Meat Price Shock and Household Inflation Expectations}
\label{tab:eq1_expectations}
\begin{threeparttable}
\footnotesize
\setlength{\tabcolsep}{3pt}
\begin{tabular}{lcccccc}
\toprule
& (1) & (2) & (3) & (4) & (5) & (6) \\
& Bivariate & Demographics & \shortstack{Province\\+ wave} & \shortstack{Region\\$\times$ wave} & \shortstack{Region $\times$ wave\\+ demographics} & \shortstack{Province +\\region $\times$ wave} \\
\midrule
Meat shock & $0.028^{**}$ & $2.492^{***}$ & $0.090^{*}$ & $0.242$ & $2.636^{***}$ & $0.119^{**}$ \\
& $(0.012)$ & $(0.812)$ & $(0.053)$ & $(0.174)$ & $(0.806)$ & $(0.050)$ \\[6pt]
Demographics &  & Yes &  &  & Yes &  \\
Province FE &  &  & Yes &  &  & Yes \\
Wave FE &  &  & Yes &  &  &  \\
Region $\times$ wave FE &  &  &  & Yes & Yes & Yes \\
\midrule
Observations & 90,133 & 22,121 & 90,133 & 90,133 & 22,121 & 90,133 \\
$R^2$ & 0.000 & 0.038 & 0.011 & 0.002 & 0.042 & 0.012 \\
\bottomrule
\end{tabular}
\begin{tablenotes}[flushleft]
\footnotesize
\item \textit{Note.} The dependent variable is the household's ordinal inflation expectation ($\mu_{it} \in \{-1, 0, 1\}$). Standard errors clustered at the province level in parentheses. Column~(6) is the preferred specification.
\item $^{*}p<0.10$, $^{**}p<0.05$, $^{***}p<0.01$.
\end{tablenotes}
\end{threeparttable}
\end{table}


Table~\ref{tab:eq1_expectations} asks whether meat-price shocks
shift expectations (equation~\eqref{eq:eq1}).
The coefficient is positive across all five columns,
indicating that higher meat-price shocks raise household
inflation expectations.

In the bivariate specification (column~1),
$\hat{\beta}_1 = 0.028$ (s.e. = 0.012, $p = 0.03$).
A one-unit increase in the cumulative meat log-change shifts
the ordinal expectation upward by 0.028 units on a
$\{-1,0,1\}$ scale.
Given that roughly 60 percent of respondents report ``up'' in
the pooled sample, the marginal effect is modest in absolute
terms but statistically clear.

Under province-plus-wave fixed effects (column~3),
$\hat{\beta}_1 = 0.090$ (s.e. = 0.053, bootstrap $p = 0.088$).
Under region-by-wave fixed effects (column~4), the point
estimate increases to $0.242$ but precision falls
(s.e. = 0.174, bootstrap $p = 0.237$).
The identifying variation for $\beta_1$ comes from
within-region cross-province differences in the meat shock.
With 31 provinces in 6 regions, region-by-wave fixed effects
absorb 18 degrees of freedom, leaving the treatment effect to
be identified from residual province-level variation within
each region-wave cell.

Column~5 adds demographics to the region-by-wave specification:
$\hat{\beta}_1 = 2.77$ (s.e. = 0.85, $p = 0.003$), again with
the reduced sample of 22,121.
The large increase in the point estimate reflects both the
sample restriction and the inclusion of controls that absorb
variance in $\mu_{ipt}$ orthogonal to the meat shock.

Why imprecision of Equation~1 does not overturn the main result.
The weaker precision under demanding fixed effects is mainly a
power issue.
The main inference continues to come from Equation~3, where the
joint object (realised inflation minus expectations) is
estimated precisely and remains negative (bootstrap
$p < 0.001$).
The sign test for overreaction does not require a cardinal
mapping of the ordinal expectation response.

% -------------------------------------------------------------
\subsection{Decomposition: Pass-Through}%
\label{sec:eq2_results}
% -------------------------------------------------------------

\begin{table}[!htbp]
\centering
\caption{Meat Price Shock and Realized Headline CPI}
\label{tab:eq2_cpi_passthrough}
\begin{threeparttable}
\begin{tabular}{lccc}
\toprule
& (1) & (2) & (3) \\
\multicolumn{4}{l}{\textit{Dep.\ var.: Realized headline CPI (\%, annualized)}} \\
\midrule
Meat shock & $-1.812^{***}$ & $-4.703^{***}$ & $-3.519^{*}$ \\
& $(0.232)$ & $(1.646)$ & $(1.954)$ \\[6pt]
Province FE &  & Yes &  \\
Wave FE &  & Yes &  \\
Region $\times$ wave FE &  &  & Yes \\
\midrule
Observations & 93 & 93 & 93 \\
$R^2$ & 0.192 & 0.704 & 0.589 \\
\bottomrule
\end{tabular}
\begin{tablenotes}[flushleft]
\footnotesize
\item \textit{Note.} Unit of observation is province $\times$ wave.
\item Standard errors clustered at the province level.
\item Dependent variable is annualized realized headline CPI inflation in the province over the subsequent inter-wave period.
\item $^{*}p<0.10$, $^{**}p<0.05$, $^{***}p<0.01$.
\end{tablenotes}
\end{threeparttable}
\end{table}


Table~\ref{tab:eq2_cpi_passthrough} estimates
equation~\eqref{eq:eq2} at the province$\times$wave level
(93 observations).
The point estimate is negative in all columns: provinces with
larger meat-price shocks experienced lower forward
headline inflation.

In the bivariate specification (column~1), $\hat{\beta}_2 =
-1.81$ (s.e. = 0.23, $p < 0.001$).
Under province-plus-wave fixed effects (column~2),
$\hat{\beta}_2 = -4.70$ (s.e. = 1.65, $p = 0.008$).
Under region-by-wave fixed effects (column~3),
$\hat{\beta}_2 = -3.52$ (s.e. = 1.95, $p = 0.08$).
The negative coefficient grows in absolute value from the
bivariate to the fixed-effect specifications, paralleling
the pattern in the forecast-error equation.

The negative $\beta_2$ reflects the ASF supply-shock reversal.
The ASF shock contracted pork supply in 2018--2019, raising
meat prices sharply.
Provinces with larger initial herd losses experienced steeper
price increases.
As herds recovered through 2020--2022, the same provinces saw
the sharpest price reversals: meat CPI fell fastest where it
had risen most.
Because meat carries a non-trivial weight in headline CPI
(roughly 3--6 percent), the reversal dragged down headline CPI
in high-shock provinces.

This reversal is central to the overreaction interpretation.
Households in high-shock provinces observed large meat-price
increases and revised their inflation expectations upward.
But the price signal was temporary: future headline inflation in
their province was not higher but if anything lower.
The household overreacted in the sense that it treated a
transitory shock as informative about future inflation when
the shock predicted the opposite.

The preferred estimate of $\hat{\beta}_2 = -3.52$ exceeds what
meat's CPI weight alone would imply ($\omega \approx 0.05$),
suggesting the meat-price reversal was correlated with declines
in other food sub-components.
This is plausible: as pork supply recovered, prices of
substitute meats (poultry, eggs) also fell, amplifying the
headline CPI decline in high-shock provinces.

% -------------------------------------------------------------
\subsection{Comparison with the Rational Benchmark}%
\label{sec:calibration}
% -------------------------------------------------------------

\begin{table}[!htbp]
\centering
\caption{Calibration: Estimated vs.\ Rational Forecast-Error Coefficient}
\label{tab:calibration_benchmark}
\begin{threeparttable}
\begin{tabular}{lccccc}
\toprule
& & \multicolumn{2}{c}{Rational bound ($\beta_3^R$)} & & \\
\cmidrule(lr){3-4}
Persistence scenario & $\rho$ & Correct belief & Worst case & $\hat{\beta}_3$ & Ratio \\
\midrule
Transitory ($\rho = 0$) & $0.00$ & $0$ & $-0.050$ & $-6.067$ & $121\times$ \\
Estimated $\hat{\rho}$ & $0.48$ & $0$ & $-0.026$ & $-6.067$ & $233\times$ \\
Moderate ($\rho = 0.5$) & $0.50$ & $0$ & $-0.025$ & $-6.067$ & $243\times$ \\
Permanent ($\rho = 1$) & $1.00$ & $0$ & $0$ & $-6.067$ & $\infty$ \\
\bottomrule
\end{tabular}
\begin{tablenotes}[flushleft]
\footnotesize
\item \textit{Note.} $\omega = 0.05$ (meat weight in headline CPI). $\hat{\rho} = 0.480$ estimated from AR(1) on monthly province-level meat CPI with province fixed effects.
\item Rational bound under correct belief: $\beta_3^R = 0$ (household knows true $\rho$).
\item Worst case: household believes shock is permanent ($\hat{\rho} = 1$), giving $\beta_3^R = -\omega(1 - \rho)$.
\item $\hat{\beta}_3 = -6.067$ from the preferred specification (region $\times$ wave FE, Table~\ref{tab:eq3_forecast_error}, column 4).
\item The ordinal expectation scale inflates the exact ratio; the point estimate should be read as an order-of-magnitude statement.
\end{tablenotes}
\end{threeparttable}
\end{table}


Table~\ref{tab:calibration_benchmark} compares the estimated
$\hat{\beta}_3$ against the rational benchmark derived in
Section~\ref{sec:rational_benchmark}.
This comparison is the paper's main quantitative check.

Under the rational benchmark, the forecast-error coefficient is
$\beta_3^R = \omega(\rho - \hat{\rho})$, where $\omega
\approx 0.05$ is meat's weight in headline CPI, $\rho$ is
the true persistence of the meat-price shock, and $\hat{\rho}$
is the household's belief about persistence.
For any $\rho$ and $\hat{\rho}$ in $[0,1]$, the rational
coefficient is bounded:
$|\beta_3^R| \leq \omega = 0.05$.
The bound is attained only in the extreme case where the
household's persistence belief is maximally wrong---for
example, the household believes the shock is permanent
($\hat{\rho} = 1$) when it is fully transitory ($\rho = 0$).

The table then calibrates the benchmark under four persistence
assumptions, corresponding to different values
of $\rho$.

Row~1 assumes the shock is fully transitory ($\rho = 0$).
If the household knows this, the rational forecast error is
zero.
In the worst case ($\hat{\rho} = 1$), the rational bound is
$-0.05$.
The estimated $\hat{\beta}_3 = -6.07$ exceeds this by a
factor of 121.

Row~2 uses the estimated persistence from an AR(1) regression
on monthly province-level meat CPI with province fixed effects:
$\hat{\rho} = 0.48$ (s.e. = 0.01, $p < 0.001$).
Under correct belief, the rational forecast error is zero.
In the worst case, $\beta_3^R = 0.05 \times (0.48 - 1) =
-0.026$.
The ratio of the estimated to the rational worst-case
coefficient is 233.

Row~3 assumes moderate persistence ($\rho = 0.5$), close to
the estimated value.
The worst-case rational bound is $-0.025$, and the ratio
is 243.

Row~4 assumes the shock is permanent ($\rho = 1$).
Under correct belief, $\beta_3^R = 0$.
Under any other belief ($\hat{\rho} < 1$), $\beta_3^R > 0$:
the household under-predicts inflation.
There is no value of $\hat{\rho}$ that generates a negative
rational forecast error when $\rho = 1$.
The ratio is infinite.

These ratios, ranging from 121 to infinity depending on the persistence
assumption, mean that the estimated forecast-error
coefficient is hard to reconcile with the rational benchmark
in Section~\ref{sec:framework}.
The gap is about economic magnitude, not statistical
significance.
The estimated coefficient is far larger than anything a rational
household could produce, even under extreme persistence
assumptions.

One might ask whether the large ratio is partly an artifact of
the ordinal expectation scale.
The \emph{sign} of the test does not depend on the mapping $k$:
whatever scaling maps the ordinal expectation into CPI units,
$\hat{\beta}_3 < 0$ rejects rational expectations.
However, the \emph{magnitude} of the ratio is sensitive to $k$.
The forecast error $\mathrm{FE}_{ipt} = \pi^{\text{headline}} -
k\,\mu_{ipt}$ inherits the scaling parameter, and a larger $k$
mechanically inflates $|\hat{\beta}_3|$ relative to the rational
bound.
The baseline mapping $k = 2$ (so that $\{-1,0,1\} \to
\{-2,0,2\}$ percentage points) yields the reported ratio;
alternative mappings in Appendix~\ref{app:scaling} shift the
ratio but do not bring it below 50.
The point estimate of $\hat{\theta}$ should therefore be
interpreted as an order-of-magnitude statement rather than a
precise calibration.

Under the diagnostic model of
Section~\ref{sec:diagnostic_extension}, the excess forecast
error identifies the product $\theta\,\omega$, where $\theta$
is the diagnostic parameter.
The implied $\hat{\theta}$ is
$(\hat{\beta}_3 - \beta_3^R) / \omega$.
Taking the most conservative rational bound ($\beta_3^R =
-0.05$), $\hat{\theta} = (-6.07 + 0.05) / 0.05 = 120.4$.
This is large.
Even dividing by 10 to account for potential measurement
amplification from the ordinal scale gives $\theta \approx 12$,
still far above the values of $\theta \in [0.5, 1.5]$
typically estimated in the professional-forecast literature
\citep{Bordalo2020b}.
The household response to a salient commodity shock is
qualitatively different from the overreaction patterns
documented in professional forecasts.

% -------------------------------------------------------------
\subsection{Visualising the Treatment Variation}%
\label{sec:event_study}
% -------------------------------------------------------------

\begin{figure}[htbp]
\centering
\includegraphics[width=0.84\textwidth]{fig_meat_cpi_spaghetti.pdf}
\caption{Province-level meat CPI trajectories, 2016--2025.
  Each line is one province; colour indicates macro-region.
  The vertical dashed line marks August 2018 (first ASF case).
  The sharp divergence in 2019--2020 followed by mean reversion
  illustrates the temporary supply-shock nature of the ASF
  episode.}
\label{fig:meat_cpi_spaghetti}
\end{figure}

Figure~\ref{fig:meat_cpi_spaghetti} displays province-level
meat CPI trajectories from 2016 to 2025.
All 31 provinces track closely before August 2018, then diverge
sharply, with some experiencing increases of 80--100 log points
and others 40--60 log points.
The trajectories converge back by 2021--2022 as herds recovered,
confirming the temporary nature of the shock.
The dispersion appears across regions, alleviating the
concern that results are driven by outliers.

\begin{figure}[htbp]
\centering
\includegraphics[width=0.80\textwidth]{fig_meat_shock_distribution.pdf}
\caption{Distribution of province-level meat-price shocks
  by CFPS wave.  Each dot is one province; boxes show
  interquartile range.  The 2020 wave captures the ASF
  peak; the 2022 wave shows the reversal.}
\label{fig:meat_shock_dist}
\end{figure}

Figure~\ref{fig:meat_shock_dist} shows the distribution of
the treatment variable by CFPS wave.
The 2020 wave has the largest and most dispersed shocks,
with an interquartile range of roughly 0.15 to 0.55 log points.
The 2022 wave shows negative shocks as meat prices reversed,
nearly uniformly across provinces.
Pre-ASF waves (2016, 2018) show modest shocks clustered
near zero, confirming that the ASF episode is the dominant
source of treatment variation.

\begin{figure}[htbp]
\centering
\includegraphics[width=0.80\textwidth]{fig_meat_vs_headline.pdf}
\caption{Meat shock vs.\ forward realised headline CPI
  at the province level.
  The fitted line shows the negative relationship: provinces
  with larger meat shocks do not experience higher subsequent
  headline inflation.}
\label{fig:meat_vs_headline}
\end{figure}

Figure~\ref{fig:meat_vs_headline} plots the meat shock against
forward headline CPI at the province level.
Each point is one province-wave cell.
The fitted line slopes downward, confirming that high-shock
provinces experienced lower subsequent headline inflation.
This provides a visual counterpart to
Table~\ref{tab:eq2_cpi_passthrough}.
The negative slope is the visual expression of the overreaction
logic: households raised expectations in response to a signal
that predicted the opposite of what they expected.

